\documentclass[11pt,a4paper]{article}
\usepackage[margin=29mm]{geometry}
\usepackage[T1]{fontenc}
\usepackage{lmodern}
\usepackage{amsmath,amssymb,amsthm,mathtools}
\usepackage{enumitem,booktabs,longtable,array}
\usepackage{xcolor,hyperref,microtype}
\usepackage{url}
\hypersetup{colorlinks=true,linkcolor=blue!45!black,citecolor=blue!45!black,urlcolor=blue!45!black}
\setlength{\emergencystretch}{2em}
\newcommand{\Q}{\mathbb Q}
\newcommand{\Z}{\mathbb Z}
\newcommand{\R}{\mathbb R}
\newcommand{\Pj}{\mathbb P}
\newcommand{\G}{\mathrm G}
\newcommand{\Sec}{\operatorname{Sec}}
\newcommand{\khat}[1]{\widehat{#1}^{\times}}
\newenvironment{statement}[3]{\par\medskip\noindent\textbf{#1. [#2] #3.}\ \itshape}{\par\medskip}
\title{Sections of the tripod over $\Q$:\\proofs, counterexamples, and unresolved steps}
\author{AI-produced by GPT-6 Astra (OpenAI, ChatGPT Pro mode)}
\date{16 November 2026}

\begin{document}
\section{Birational and valuative results}
% Birational agent fragment.  Insert in Part B.  Bibliography: BIR_* keys.
% All section maps below are continuous; geometric conjugacy is specified.
\subsection{The birational comparison and its local consequences}

Fix compatible geometric generic base points.  For a smooth geometrically
connected curve $X/k$, put $K=k(X)$ and write
\[
 q_X:G_K\twoheadrightarrow\pi_1(X),\qquad
 q_{X,\mathrm{geom}}:G_{K\bar k}\twoheadrightarrow\pi_1(X_{\bar k}).
\]
These are continuous surjections over $G_k$.  They induce a map on
continuous section classes, with conjugation by the respective geometric
kernels.  Its image is denoted $\Sec^{\mathrm{bir}}(X/k)$.

\begin{statement}{B.1}{Known}{The precise local and adelic consequences of a birational lift}
Let $k$ be a number field and $X/k$ a hyperbolic curve with smooth
completion $\overline X$.  A section in $\Sec^{\mathrm{bir}}(X/k)$ is
locally a point or cusp section.  At finite places this supplies centers
$x_v\in\overline X(k_v)$; at real places the section specifies a connected
component.  Such a section is either globally cuspidal or is integral
on $X$ at a set of finite places of Dirichlet density one.

A birationally adelic section means one admitting a single birational
lift whose image on every nonempty open $V\subseteq X$ is adelic on $V$
or cuspidal.  Here adelic requires ordinary local sections everywhere
and integrality at all but finitely many finite places.  Over a totally
real or imaginary quadratic number field, every noncuspidal
birationally adelic section is a rational point section.  In particular,
every adelic section of the tripod over $\Q$ is a rational point section.
\end{statement}
\begin{proof}
These are respectively Proposition 1, Theorem B, the definition in
\S2.2.3, Theorem A, and Corollary 6 of
\cite{BIR_Stix2015}; the last applies because $t$ is a nonconstant
unit on the tripod.  The comparison maps between the hypotheses are
the actual inclusions
\[
 \Sec^{\mathrm{ba}}(X/k)\lhook\joinrel\longrightarrow
 \Sec^{\mathrm{bir}}(X/k)\lhook\joinrel\longrightarrow\Sec(X/k).
\]
The first forgets the adelic conditions on a chosen lift, and the second
forgets liftability.  The density assertion does not prove the cofinite
integrality required in the definition of the first set.

For clarity we specify the local comparison underlying Proposition 1.
For each finite section-neighbourhood, Koenigsmann's Corollary 2.6
supplies a point on its proper completion at every completion of $k$.
The spaces of such points are compact.  A common refinement is again
a section-neighbourhood, so the closed compatibility conditions have
the finite intersection property.  Their inverse limit is a compatible
local pro-point, and its images induce the local point or cusp section.
This is the compactness argument of \cite[Lemma 1.7]{BIR_Koenigsmann2005}.
It does not identify
$G_{k_v(X)}\to G_{k(X)}\times_{G_k}G_{k_v}$ with an isomorphism.

The full local birational assertion over finite extensions of $\Q_p$
is \cite[Proposition 2.4(2), Definition 1.4]{BIR_Koenigsmann2005}.
Theorem A of \cite{BIR_Pop2010} gives its specified elementary-abelian
$p$-quotient refinement with the additional assumption $\mu_p\subseteq k$.
That theorem uses the maximal elementary abelian $p$ extension and
then the same extension of that field, with the stated lift between
them.  We make no identification with a lower-central class-two
quotient.  No October 2026 local section theorem is used here.
\end{proof}

\begin{statement}{B.2}{Proved}{Actual birational lifts of the point and cusp loci}
For $U=\Pj^1_\Q\setminus\{0,1,\infty\}$, every rational point section
and every section in a rational cusp packet has a continuous lift
through $q_U:G_{\Q(t)}\twoheadrightarrow\pi_1(U)$.  This holds for
chosen representatives as well as for geometric conjugacy classes.
\end{statement}
\begin{proof}
At $x\in\Pj^1(\Q)$ take $z=t-x$, or $z=1/t$ at infinity.  A chosen
prolongation of the $z$-adic valuation gives
\[
 1\longrightarrow I_x\simeq\widehat\Z(1)
 \longrightarrow D_x\longrightarrow G_\Q\longrightarrow1.
\]
We split this sequence by fixing compatible roots of $z$ and acting
on the coefficients of algebraic Puiseux series.  A finite extension
of $\Q((z))$ embeds in $L((z^{1/n}))$ for a finite extension $L/\Q$
and an integer $n$: the residue characteristic is zero, so all
ramification is tame, and Hensel's lemma extracts roots of units.
The coefficientwise actions are compatible and continuous at every
finite level.  Restricting to the separable closure of the henselization
of $\Q(t)$ therefore gives the asserted continuous splitting in $D_x$.

For $x\in U(\Q)$, $q_U$ kills $I_x$ and sends that splitting to
the point section $s_x$.  For a cusp, the map
$D_x\to D_x^U$ is an isomorphism: it induces the identity on $G_\Q$
and an isomorphism on $\widehat\Z(1)$.  The latter inertia is injective
in the geometric free profinite group of the tripod; its peripheral
element is primitive.  Consequently any splitting in the cusp
decomposition group, including every profinite tangential parameter,
lifts through this isomorphism.  Finally, a geometric conjugator
in $\pi_1(U_{\bar\Q})$ lifts through $q_{U,\mathrm{geom}}$.
Conjugating the lifted section proves the assertion for representatives.
\end{proof}

\begin{statement}{B.3}{Proved}{Finite-\'etale transfer of the point and cusp loci}
Let $f:Y\to X$ be a finite \'etale map of smooth geometrically connected
hyperbolic curves over a characteristic-zero field $k$.  The full
inverse image of the point-or-cusp locus under
\[
 f_*:\Sec(Y/k)\longrightarrow\Sec(X/k)
\]
is the point-or-cusp locus of $Y$.  Above the class of $s_x$, for
$x\in X(k)$, it is precisely the set of point classes belonging to
$Y_x(k)$.
\end{statement}
\begin{proof}
Compatible base points identify $\pi_1(Y)$ with an open subgroup
$H\subset\pi_1(X)$, and its geometric subgroup with
$H\cap\pi_1(X_{\bar k})$.  A geometric conjugation of a section
replaces $H$ by the corresponding conjugate subgroup, so it suffices
to work with actual inclusions of section images.

The inclusion $s_x(G_k)\subset H$ says exactly that the associated
sheet of the finite fibre $Y_x(\bar k)$ is fixed by $G_k$.  It is
therefore a point $y\in Y_x(k)$.  With the fibre path at $y$, the
sections $r$ and $s_y$ have the same image $s_x$ in $\pi_1(X)$.
Injectivity of $f_*$ gives their equality.

For a cusp $c$, normalize $\overline X$ in $k(Y)$.  For the chosen
prolongation of $c$, the subgroup $H\cap D_c$ is the decomposition
group at the point $d$ above $c$ on that normalization.  Its image
in $G_k$ is $G_{k(d)}$.  Containing the image of a section forces
this image to be all of $G_k$, and hence $k(d)=k$.  The lifted
section thus lies in the rational cusp packet at $d$.  Conversely,
point and cusp sections map to their corresponding point and cusp
packets by the same fibre and valuation functors.

In particular a Belyi morphism $C\to\Pj^1$ applies this argument to
the finite \'etale open $C\setminus\beta^{-1}\{0,1,\infty\}\to U$.
A section of the proper curve $C$ must first lift to that open to
enter the displayed map.  The transfer assertion itself does not
supply such a lift.
\end{proof}

\subsection{The full group, its relative central quotient, and the omitted kernel}

\begin{statement}{B.4}{Proved}{A compatible splitting of every relative central puncture quotient}
Let $k$ have characteristic zero, let
$U=\Pj^1_k\setminus\{0,1,\infty\}$, and let $D\subset U$ be finite
and reduced.  Put $W=U\setminus D$ and
$T_D=\operatorname{Res}_{D/k}\mathbb G_m$.  For
\[
 N=\ker(\pi_1(W)\to\pi_1(U)),\qquad
 N_{\mathrm{cc}}=\overline{[N,\pi_1(W_{\bar k})]},
\]
the morphism $(j,t-a):W\to U\times_kT_D$, where $a$ is the
tautological coordinate of $D$, induces an isomorphism over $G_k$
\[
 \pi_1(W)/N_{\mathrm{cc}}
 \xrightarrow{\ \sim\ }
 \pi_1(U)\times_{G_k}\pi_1(T_D).                 \tag{B.4.1}
\]
The projection to $\pi_1(U)$ has a continuous splitting, and these
splittings are compatible under deletion of further finite divisors.
\end{statement}
\begin{proof}
The function $t-a$ is invertible on $W\times_kD$, so restriction of
scalars gives the stated $k$-morphism.  Over $\bar k$, enumerate
$D(\bar k)=\{a_1,\dots,a_m\}$.  The free profinite group of $W$
has free generators $\gamma_0,\gamma_1,\delta_1,\dots,\delta_m$.
The kernel $N$ is normally generated by the $\delta_i$; imposing
$[N,\pi_1(W_{\bar k})]=1$ makes precisely these generators central.
The resulting universal profinite group is
$\widehat F_2\times\widehat\Z^m$.
The winding numbers of $t-a_i$ are zero at $0,1$ and are
$\delta_{ij}$ at $a_j$.  Thus the displayed morphism is an isomorphism
on geometric kernels.  It induces the identity on $G_k$, proving
(B.4.1) on arithmetic groups.  This also verifies explicitly, in this
case, the relative cuspidally central comparison of
\cite[Definition 2.1 and Propositions 3.1, 3.4, 3.5]{BIR_BorneEmsalemStix2015}.
Those numbers are from the attached arXiv v2; the opened earlier
author PDF numbers them Definition 1 and Propositions 8, 11, 12.

The identity $1\in T_D(k)$ induces $s_1:G_k\to\pi_1(T_D)$.
Under (B.4.1) the required splitting is the continuous homomorphism
\[
 e\longmapsto(e,s_1(p_U(e))),\qquad e\in\pi_1(U).
\]
For a disjoint union of divisors the torus is the product and its
identity is the tuple of identities.  Forgetting a divisor is exactly
the corresponding projection, proving compatibility.
The remaining map to the full group is
$\pi_1(W)\twoheadrightarrow\pi_1(W)/N_{\mathrm{cc}}$; its kernel
is the specified group $N_{\mathrm{cc}}$.
\end{proof}

\begin{statement}{B.5}{Refuted}{A chosen central lift need not have a birational lift}
For $k=\Q$, $D=\{3\}$ and $W=U\setminus\{3\}$, the central
section $(s_2,s_1)$ in (B.4.1) is not the image of any section of
$G_{\Q(t)}\to G_\Q$.  Its projection $s_2$ does have birational
lifts.
\end{statement}
\begin{proof}
Suppose $\tau$ were such a birational lift.  Its section on $W$ is
locally a point or cusp by B.1.  At $3$ its image on $U$ is $s_2$.
An ordinary local center $x$ satisfies
$\kappa(x)=\kappa(2)$ in $\widehat{\Q_3^\times}$ and therefore
$x=2$.  Indeed
$\Q_3^\times\simeq3^\Z\times\mu_2\times(1+3\Z_3)$ injects
into its power completion, as follows on the last factor from the
$3$-adic logarithm.

The three cusp possibilities are also excluded by explicit units.
At $0$ the coordinate $1-t$ has Kummer value $1$, whereas $s_2$
gives $-1$.  At $1$ the coordinate $t$ has value $1$, whereas
$s_2$ gives $2$.  At infinity the unit $(1-t)/t$ has value $-1$,
whereas $s_2$ gives $-1/2$.  Each pair differs already in
$\Q_3^\times/(\Q_3^\times)^2$, since $-1$ and $2$ are nonsquares.
The local center on $W$ is consequently $2$ as well; a cusp at
the newly deleted point $3$ would map to the ordinary point $3$
on $U$.  The coordinate $t-3$ then has local Kummer class $-1$.
The prescribed central section gives class $1$, a contradiction
modulo squares.  The finite residue calculations are recorded in
\texttt{birational\_central\_certificate.py} and its output.

B.2 provides birational lifts of $s_2$; their $t-3$ coordinate is
$-1$.  The counterexample concerns the chosen central lift.  In
particular, the reciprocity criterion proved for canonical central
lifts in AP D.27 addresses the middle group in
\[
 \pi_1(W)\twoheadrightarrow\pi_1(W)/N_{\mathrm{cc}}
       \twoheadrightarrow\pi_1(U).
\]
Its proof expressly retains the nonabelian lifting problem.  Promoting
that criterion to a full lift, as in the wording of R9/A8, requires
a further proof through the first arrow.  This counterexample does
not assert that the canonical criterion is insufficient for every
possible choice of its extra data.
\end{proof}

\begin{statement}{B.6}{Proved}{The configuration-space interpretation of a puncture lift}
Let $F=(U\times_\Q U)\setminus\Delta$ with coordinates $(t,z)$ and
projections $p_1,p_2$.  For a chosen section $s$ of $U$ put
\[
 \Pi_s=\pi_1(F)\times_{\pi_1(U),p_1,s}G_\Q,
 \qquad h_s:\Pi_s\longrightarrow\pi_1(U),\quad h_s=p_2.
\]
This is a continuous extension of $G_\Q$ by $\widehat F_3$.
For $x\in U(\Q)$, full lifts of $s$ to $\pi_1(U\setminus\{x\})$
are exactly sections $r:G_\Q\to\Pi_s$ satisfying
$h_sr=s_x$, with compatible representatives and base paths.
\end{statement}
\begin{proof}
The geometric fibre of either projection is a four-punctured projective
line.  The topological fibre sequence over $\C$, the vanishing of
$\pi_2$ of the base punctured curve, and the finite-cover comparison
give the asserted exact geometric fundamental-group sequence;
characteristic-zero base change and the arithmetic sequences give
the sequence over $\Q$.  The map
\[
 i_x:U\setminus\{x\}\longrightarrow F,\qquad t\longmapsto(t,x)
\]
induces an isomorphism
\[
 \pi_1(U\setminus\{x\})
 \xrightarrow{\ \sim\ }p_2^{-1}(s_x(G_\Q)).       \tag{B.6.1}
\]
The fibre-sequence diagram induces the identity on the geometric
fibre $\widehat F_3$ and on $G_\Q$, which proves (B.6.1).
Taking the pullback along $s$ now proves both directions of the claim.

There is always a tangential splitting of $\Pi_s$ at the relative
mark $z=0$.  Its second projection is a cusp section at $0$, so it
does not furnish the prescribed ordinary projection $s_x$ in
(B.6.1).  The double cover
$y^2=z(z-1)(z-t)$ over $F$ is \'etale, and its compactified fibres
are the Legendre elliptic curves; hence its fibre cohomology gives
the actual Legendre representation after pullback along $s$.
These maps retain the interpolation requirement in
\cite[Theorem C]{BIR_Bresciani2024}; individual central splittings
in B.4 have not supplied a section with the projection in (B.6.1).
\end{proof}

\begin{statement}{B.7}{Proved}{The diagonal inertia omitted by transcendental birational base change}
Let $x,T$ be independent over $\Q$, fix
$\Omega=\Q(x,T)^{\mathrm{sep}}$, and let $L_x,L_T$ be the algebraic
closures of $\Q(x),\Q(T)$ in $\Omega$.  Set $L=L_xL_T$.
Restriction gives a continuous exact sequence
\[
 1\longrightarrow G_L\longrightarrow G_{\Q(x,T)}
 \xrightarrow{R}G_{\Q(x)}\times_{G_\Q}G_{\Q(T)}
 \longrightarrow1.                              \tag{B.7.1}
\]
For a prolongation of the divisor $x-T=0$, its nontrivial inertia
$I_\Delta\simeq\widehat\Z(1)$ is contained in $G_L$.
\end{statement}
\begin{proof}
The common algebraic constants form $\bar\Q$.  A common element
$\alpha\in L_x\cap L_T$ transcendental over $\Q$ would make both
$x$ and $T$ algebraic over $\Q(\alpha)$, contradicting their
independence.  The stronger linear disjointness needed here follows
as follows.  Finite extensions of $\bar\Q(x)$ and $\bar\Q(T)$
are function fields of geometrically integral curves.  Their product
is integral, with independent coordinates $x,T$.  Thus the two
fields inject into their tensor compositum without additional
relations.  Passing to their unions proves that $L_x,L_T$ are
linearly disjoint over $\bar\Q$.  Automorphisms agreeing on the
constants consequently extend uniquely to their compositum, and
then extend to $\Omega$.  This proves the surjectivity in (B.7.1)
and its stated kernel.

The valuation $v_\Delta$ is trivial on each of $\Q(x)^\times$
and $\Q(T)^\times$.  A separable primitive polynomial for a finite
extension of either one-variable field has integral coefficients
and unit discriminant for $v_\Delta$.  The extension it defines
is therefore unramified at that valuation.  Finite composita remain
unramified, by the finite \'etale category of the henselian valuation
ring.  Thus the inertia of $v_\Delta$ acts trivially on $L$.
Residue characteristic zero gives $I_\Delta\simeq\widehat\Z(1)$;
the extensions adjoining $(x-T)^{1/n}$ show its nontrivial finite
quotients and are absent from $L$ for $n>1$.

More explicitly, a chosen birational section
$\tau:G_\Q\to G_{\Q(x)}$ defines
\[
 b_\tau:G_{\Q(T)}\longrightarrow
          G_{\Q(x)}\times_{G_\Q}G_{\Q(T)},\qquad
 g\longmapsto(\tau(g|_{\bar\Q}),g).
\]
The continuous pullback $E_\tau$ of (B.7.1) along $b_\tau$ has
exact sequence
\[
 1\longrightarrow G_L\longrightarrow E_\tau
       \longrightarrow G_{\Q(T)}\longrightarrow1. \tag{B.7.2}
\]
A lift preserving these two restrictions is precisely a splitting
of (B.7.2), by the universal property of the pullback.  Nontriviality
of its kernel proves no nonsplitting assertion: the sections in B.2
provide examples where it splits.

For a model section $s$, geometric base change instead gives
$\pi_1(U_{\Q(T)})\simeq\pi_1(U)\times_{G_\Q}G_{\Q(T)}$.
Its $t$-birational lifting problem is the pullback of
$G_{\Q(x,T)}\to\pi_1(U_{\Q(T)})$ along $s_{\Q(T)}$.
The map from (B.7.2) to that pullback forgets the chosen restriction
$\tau$ on $G_{\Q(x)}$.  No splitting for an arbitrary $s$ is
constructed by this map.  The additional $\Q(T)$ hypothesis in
\cite[Theorem A]{BIR_Bresciani2024} is therefore retained in its
application here.  The definitions and Theorems 1.3/1.7 in
\cite{BIR_Pop2026} likewise retain a transcendental-base lift;
we use none of that preprint's conclusions.
\end{proof}

\subsection{A single coefficient system with all local elliptic realizations}

\begin{statement}{B.8}{Proved}{An infinitely ramified integral representation locally realized by elliptic curves}
There is a continuous representation
$\rho:G_\Q\to\mathrm{GL}_2(\Z_7)$ with surjective residual
representation, determinant $\chi_7$, and infinitely many ramified
primes, such that for every finite prime $q$ there are an elliptic
curve $E_q/\Q_q$ and an isomorphism
\[
 \rho|_{G_{\Q_q}}\simeq T_7(E_q)
 \quad\text{of }\Z_7[G_{\Q_q}]\text{-modules}.      \tag{B.8.1}
\]
One can take $E_7$ to be $E:y^2=x(x-1)(x-3)$.  The representation
is absolutely irreducible; its ramification set has density zero;
at every unramified $q\ne7$ its polynomial is
$X^2-a_qX+q$ with $a_q\in\Z$ and $|a_q|\le2\sqrt q$.
\end{statement}
\begin{proof}
The displayed curve has $\Delta=576$ and $c_4=112$.  At $3$ it
has multiplicative reduction and minimal discriminant valuation $2$.
The equation is minimal there since its discriminant valuation is
less than $12$.  Tate uniformization gives a nontrivial transvection
on $E[7]$ from inertia, because $7\nmid2$.  In the nonsplit case
the required quadratic twist is unramified and preserves this inertia
calculation.  Also $\#E(\mathbf F_5)=4$, so $a_5=2$ and the
discriminant of $X^2-2X+5$ modulo $7$ is the nonsquare $5$.
Thus the residual representation is irreducible over $\mathbf F_7$.

To check the full residual image, let $u$ be the transvection and
$M$ its fixed line.  Irreducibility gives $gM\ne M$.  A basis on
$M,gM$ makes $u,gug^{-1}$ nontrivial upper and lower unipotents.
Their powers give every upper and lower elementary matrix over the
prime field $\mathbf F_7$ and therefore generate
$\mathrm{SL}_2(\mathbf F_7)$.  The determinant is cyclotomic and
surjective, so the image is $\mathrm{GL}_2(\mathbf F_7)$.
This also proves absolute irreducibility.  The exact point count
and invariant calculations are in
\texttt{locally\_elliptic\_residual\_certificate.py} and its output.

We now use the construction, not merely the statement, of
\cite[\S2.2, Main Theorem and its following remark]{BIR_KLR2005}.
Set $p=7$, $k=1$, $\bar\rho=E[7]$, and fix determinant $\chi_7$.
Include $2,3,7$ in the initial finite set $S$ and enlarge it as in
their Lemma 7.  Choose $E[49]$ as the initial lift and $T_7(E)$
as every prescribed characteristic-zero local lift on $S$.
For their finite initial set $R_2$, choose the Frobenius polynomials
of $E$; these satisfy the initial congruences and are pure.  The
subsequent trace choices in their construction lie in the weight-one
Hasse interval.  At each new auxiliary prime choose the branch
$A=0$ of their equation (3) and a nonzero characteristic-zero value
of $B$ lifting its finite-level value.  The induction freezes these
chosen local lifts.  The cited construction permits infinitely many
auxiliary primes and forces them to ramify.  This supplies $\rho$,
with no initial purity exceptions.  Its absolute irreducibility
follows from its absolutely irreducible residual representation;
its density assertion is \cite[Theorem 1]{BIR_KhareRajan2001}.

We verify (B.8.1) integrally at every place.  At a new auxiliary
prime $q$, the nice-prime condition gives $q\not\equiv\pm1\pmod7$.
The fixed local representation of equation (3), with $A=0$, has
\[
 \rho(\sigma_q)=\psi(\sigma_q)
          \begin{pmatrix}q&0\\0&1\end{pmatrix},\qquad
 \rho(\tau_q)=\begin{pmatrix}1&B\\0&1\end{pmatrix},
 \quad 0\ne B\in\Z_7.                           \tag{B.8.2}
\]
The twist $\psi$ is unramified and quadratic: residual inertia is
trivial, the deformation has the displayed unipotent inertia, and
the fixed determinant forces $\psi^2=1$.  Distinct residual
Frobenius eigenvalues justify integral diagonalization.  Write
$B=7^m b$ with $b\in\Z_7^\times$; conjugation by
$\operatorname{diag}(b^{-1},1)$ changes $B$ to $7^m$ and fixes
the diagonal Frobenius.  The Tate curve with parameter $q^{7^m}$
has exactly this integral representation before the twist.
Its $7^n$-torsion is generated by $\zeta_{7^n}$ and
$(q^{7^m})^{1/7^n}$; choose a Frobenius lift fixing compatible
roots of $q$.  Its matrices are those just obtained.  The
uniformization and its torsion exact sequence are
\cite[Theorem 1 and the discussion of {$E_t[m]$}, pp.~18--19]{BIR_Tate1995}.
Twisting the curve by $\psi$ proves (B.8.1) at $q$, with the lattice
unchanged except for the stated matrix in $\mathrm{GL}_2(\Z_7)$.

At an unramified prime outside $S$ we have $q\ge5$, $q\ne7$.
Every integer in its Hasse interval occurs as an elliptic trace over
$\mathbf F_q$ by \cite[Theorem 4.1]{BIR_Waterhouse1969}:
nonzero such integers are prime to $q$, and zero is the allowed
odd-exponent case.  Lift a corresponding Weierstrass equation to
$\Z_q$.  The resulting $E'_q/\Q_q$ has good reduction.
Since $q$ is prime, $a_q^2-4q\ne0$; thus both unramified rational
representations are semisimple and are isomorphic by their
Frobenius polynomials.

Identify their rational Tate spaces and let $L$ be the stable
lattice furnished by $\rho$.  After a scalar homothety arrange
\[
 T_7(E'_q)\subset L\subset7^{-n}T_7(E'_q).
\]
The quotient $H=L/T_7(E'_q)$ is a $G_{\Q_q}$-stable finite
subgroup of $E'_q[7^n]$.  The quotient elliptic curve $E_q=E'_q/H$
is defined over $\Q_q$.  The rational identification induced by
the isogeny identifies its integral Tate lattice with $L$:
the quotient of the larger lattice by the smaller is precisely
the kernel of that isogeny in the divisible Tate module.
The subgroup extends to a finite \'etale subgroup of the good
integral model, since $7$ is invertible in $\Z_q$; the quotient
again has good reduction.  Scalar homothety preserves the
isomorphism class of the integral Galois module.  This establishes
the required integral match, without selecting a preferred
Frobenius lattice in its isogeny class.

At every place in the original $S$, including $2,3,7$, the local
representation was fixed to $T_7(E)$.  In particular $E$ has good
reduction at $7$, because $7\nmid\Delta$.  This covers every
finite place, including the coefficient prime.
\end{proof}

The exact comparison with the full Legendre data is the forgetful map
\[
 \bigl((\rho_\ell)_\ell,(E_q)_q,(\iota_{q,\ell})_{q,\ell}\bigr)
 \longmapsto\bigl(\rho_7,(E_q)_q,(\iota_{q,7})_q\bigr),          \tag{B.8.3}
\]
where $\iota_{q,\ell}$ are integral local identifications with a
single common $E_q$.  B.8 gives an infinitely ramified object in
the target.  We construct no inverse image under (B.8.3), or under
the further map from actual sections by Legendre monodromy.
Thus this result addresses exactly what can be inferred from one
coefficient, including all of its local elliptic realizations.

\subsection{Valuations and the local-to-global inertia map}

\begin{statement}{B.9}{Proved}{The covering-germ criterion with its exact quantifiers}
For a continuous homomorphism $i:H\to G$ of profinite groups the
following are equivalent: $i$ is injective; every finite transitive
continuous $H$-set is an orbit in the restriction of a finite
continuous $G$-set; every such $H$-set is a quotient of such an orbit.
\end{statement}
\begin{proof}
For an injection, compactness identifies $H$ with a closed subgroup.
Given an open subgroup $V\subset H$, choose an open normal
$N\triangleleft G$ with $N\cap H\subset V$.  The subgroup $VN$
is open and $H\cap VN=V$, so the orbit of $VN$ in $G/VN$ is
exactly $H/V$.  The second assertion implies the third.  Conversely,
$\ker i$ acts trivially on all restricted $G$-sets and their
quotients.  The third assertion therefore places it in every open
normal subgroup of $H$, whose intersection is trivial.

For a geometric strict henselian normal-crossings neighbourhood,
the finite-cover correspondence identifies orbits with connected
components of restricted covers.  In characteristic zero its local
group is $H=\widehat\Z(1)^{I(x)}$: adjoining all roots of the
separated boundary parameters removes tame ramification, and the
remaining \'etale covers of a strictly henselian local scheme are
trivial.  The criterion is consequently a criterion for the
injectivity of this actual local-to-global homomorphism.  This
reproves the profinite assertion of AP G.1 and supplies the precise
meaning of ``sub-germ'' in cote 149, scan 130.  It does not assert
that the homomorphism is always injective.
\end{proof}

\begin{statement}{B.10}{Refuted}{Distinct rational-residue valuations need not have distinct nonzero inertia images}
Let $C/\Q$ be the smooth projective curve $y^2=x^5-x$, and put
\[
 V=(\Pj^1_\Q\setminus\{0,1,\infty\})\times C.
\]
This is a strict hyperbolic polycurve.  There are distinct rank-two
valuation rings of $\Q(V)$ with residue field $\Q$ whose inertia
images in $\pi_1(V_{\bar\Q})$ are the same nontrivial subgroup,
after compatible choices of paths.  Intrinsically their conjugacy
classes of subgroups are equal.
\end{statement}
\begin{proof}
The polynomial $x^5-x$ has five distinct roots and the double cover
is also branched at infinity, so $C$ has genus two.  It has rational
points $c_1=(0,0)$ and $c_2=(1,0)$.  Projection $V\to C$ has
tripod fibres, giving the asserted hyperbolic tower.

Let $t$ denote the first coordinate.  Compose the divisorial
valuation at $t=0$ with the point valuation at $c_i$ on its residue
field $\Q(C)$.  Explicitly, with lexicographic order,
\[
 v_i(f)=\left(\operatorname{ord}_{t=0}f,
       \operatorname{ord}_{c_i}
       \overline{t^{-\operatorname{ord}_{t=0}f}f}\right)
       \in\Z^2.                                  \tag{B.10.1}
\]
Its final residue field is $\Q$.  These valuation rings are distinct:
the function $x$ has $v_1(x)=(0,2)$ and $v_2(x)=(0,0)$.
The first order follows because $y$ is a parameter at $c_1$ and
the equation gives a simple branch in $x$.

The geometric product comparison gives
\[
 \pi_1(V_{\bar\Q})\simeq
       \widehat F_2\times\pi_1(C_{\bar\Q}).
\]
In the first projection, inertia of either composite valuation
is contained in cusp inertia $I_0$.  Its second projection is
trivial: every finite \'etale cover of $C$ is unramified at the
interior point $c_i$.  Finally the first-stage divisorial inertia
at $t=0$ is contained in the inertia of the composite valuation
and maps onto $I_0$.  Hence for both $i$
\[
                   \Delta_{v_i}=I_0\times1\ne1.   \tag{B.10.2}
\]
The containment of first-stage inertia follows also directly from
the tower of decomposition groups: it acts trivially on the entire
algebraic residue field and therefore preserves the second valuation
and acts trivially on its final residue.  Paths in the $C$ factor
do not change $I_0\times1$, proving the path-independent conjugacy
assertion.  This refutes the universal distinctness proposed in
cote 149, scan 128, even within strict hyperbolic polycurves.
\end{proof}

\begin{statement}{B.11}{Refuted}{The unrestricted converse to the center criterion}
There is a smooth geometrically connected surface $W/\Q$ and a
rank-two valuation $v$ of $\Q(W)$ with residue field $\Q$, with
no center on $W$, whose full inertia image in $\pi_1(W)$ is trivial.
\end{statement}
\begin{proof}
Use $C,c_1$ from B.10, set $Y=C\times C$, take
$P=(c_1,c_1)$, and put $W=Y\setminus\{P\}$.
Purity at the regular codimension-two point gives an equivalence
between the finite \'etale categories of $Y$ and $W$, hence
$\pi_1(W)\simeq\pi_1(Y)$; the precise local extension theorem is
\cite[Lemma 58.21.3, Tag 0BJE]{BIR_Stacks}.
Blow up $P$ and choose a rational point $e$ of the exceptional
curve $E\simeq\Pj^1_\Q$.  Compose the divisorial valuation of $E$
with the point valuation at $e$.  This gives a rank-two valuation
with residue $\Q$, whose center on $Y$ is $P$.
Separatedness makes that center unique, so it has no center on $W$.

Every finite \'etale cover of $W$ extends to $Y$, then pulls back
\'etale to its blowup.  The resulting cover of the valuation ring
is unramified at both stages of this flag.  Thus inertia acts
trivially on every finite quotient of $\pi_1(W)$ and its full
image is trivial.  This answers the unrestricted converse in
cote 149, scan 120, negatively.  We supply no hyperbolic-polycurve
tower for $W$; the stronger margin speculation imposing elementary
hyperbolic factors is not refuted by this example.
\end{proof}

\begin{statement}{B.12}{Proved}{Explicit blowup and minimal-surface inertia maps}
For the strict polycurve $V$ of B.10, a blowup of the standard
compactification $X=\Pj^1\times C$ at $(0,c_1)$ has a boundary
crossing with local-to-global map
\[
 \widehat\Z(1)^2\longrightarrow I_0\times1,
                 (r,s)\longmapsto r+s.            \tag{B.12.1}
\]
The standard compactification itself has injective boundary inertia
and has no exceptional curve of the first kind.  For arbitrary smooth
opens, even this last surface minimality condition does not ensure
injectivity: on $\Pj^2$ with boundary $xy=0$ the map at $[0:0:1]$ is
\[
 \widehat\Z(1)^2\longrightarrow\widehat\Z(1),
                    (r,s)\longmapsto r-s.         \tag{B.12.2}
\]
\end{statement}
\begin{proof}
At $(0,c_1)$ use parameters $t,z$, with $z$ a parameter of $C$.
In the blowup chart $t=az$, the strict boundary and exceptional
curve are $a=0$ and $z=0$.  Both loops wind once in $t$, while
the $z$-loop is killed in $\pi_1(C)$ because its center is an
interior point.  This gives (B.12.1), whose kernel is
$\{(r,-r)\}$.  In $X$ the three boundary components are disjoint
and each has its injective single cusp inertia.  Every rational
curve in $X$ maps constantly to the genus-two curve $C$, by
Riemann--Hurwitz; it is a fibre and has self-intersection zero.
Thus $X$ has no $(-1)$-curve.

For (B.12.2), the complement is $\mathbb G_m\times\mathbb A^1$
with coordinates $u=x/y$, $v=z/y$.  In the chart $z=1$ the
orders of $u$ on the two boundary branches are $(1,-1)$.
Kummer characters therefore give the displayed map and its diagonal
kernel.  The projective plane has no $(-1)$-curve: every curve has
self-intersection equal to the square of its positive degree.
This reproves AP G.2 in a surface which is minimal in that precise
sense.  The complement has an affine-line factor, so the calculation
does not settle an existential assertion about minimal
compactifications of strict hyperbolic polycurves.  The term
``minimal'' in cote 149, scan 130, requires that further specification.
The finite reductions of (B.12.1)--(B.12.2) are checked in
\texttt{birational\_valuation\_certificate.py} and its output.
\end{proof}

\begin{statement}{B.13}{Proved}{The actual comparison from a birational section to Kummer local data}
A continuous birational section $\tau:G_\Q\to G_{\Q(t)}$ gives
classes $u_c\in\widehat{\Q^\times}$ for every $c\in\Q$ and
compatible local centers $x_v\in\Pj^1(\Q_v)$, with a choice of
real representative, such that
\[
 u_c|_v=\kappa_v(x_v-c)
       \quad\text{whenever }x_v\ne c,\infty.      \tag{B.13.1}
\]
This is a well-defined comparison on geometric conjugacy classes,
apart from the stated choice of real representative.
For $r\in\Pj^1(\Q)$, its inverse image over the tuple with all
finite centers equal to $r$ consists exactly of the birational
lifts whose model section on $U$ is $s_r$ when $r\in U(\Q)$,
or lies in the cusp packet at $r$ when $r\in\{0,1,\infty\}$.
At the real place only the corresponding component is intrinsic.
\end{statement}
\begin{proof}
Define $u_c=\tau^*\kappa(t-c)$ using the continuous Kummer class
of the rational function.  Geometric conjugation acts trivially
on this pulled-back cohomology class.  Apply the compact inverse
limit in B.1 to all finite birational section-neighbourhoods;
at the real place choose one point in that compatible inverse
limit, rather than a separately chosen representative at each
level.  For a modulus $n$, represent $u_c$ by $a_n\in\Q^\times$.
The twisted Kummer cover
\[
                       z^n=(t-c)/a_n
\]
has a $\tau$-fixed generic sheet because its pullback class is
zero.  The compatible local pro-point on its proper completion,
above $x_v\ne c,\infty$, lies in the \'etale fibre and gives
$(x_v-c)/a_n\in(\Q_v^\times)^n$.  All moduli yield (B.13.1).

Thus the proposed configurations in A12 receive an actual map
from birational sections.  The counterexamples to the proposed
rigidity of all configurations elsewhere in this paper do not
construct inverse images under this map.  We infer no rational
place from (B.13.1) alone.

For the inverse-image assertion, pass to a section-neighbourhood
whose smooth completion has positive genus.  All its finite
local centers lie above $r$, in a finite closed subscheme.
Stix's Theorem 4 (Stoll's finite-support theorem), followed by
the inverse-limit argument in the proof of
\cite[Corollary 5]{BIR_Stix2015}, supplies a compatible rational
pro-point above $r$.  Indeed the same finite subscheme condition
holds in every further neighbourhood, and the resulting nonempty
finite sets of rational points have nonempty inverse limit.
Thus its model section is a rational point section or is cuspidal,
and the center forces that point, or cusp label, to be $r$.
Conversely those point and cusp sections have the indicated
centers at every finite place, and every birational lift has the
same local model sections.  This proves both inclusions without
replacing a completion's function-field Galois group by a fibre
product.  Literal (G) is false; this inverse-image statement does
not use it as a rigidity premise.
\end{proof}

\begin{statement}{B.14}{Proved}{A full lift over a genuinely bad marked base cusp}
For $F=M_{0,5}=(U\times U)\setminus\Delta\to U_t$, there is a
section whose projected section is cuspidal at $t=0$, whose marked
generic fibre there has no potentially good marked reduction,
and which has a full birational lift.  The lift has a compatible
invariant valuation above every one-variable subfield of $\Q(t,z)$.
\end{statement}
\begin{proof}
In $\Q(t,z)$ take the rank-two flag valuation given first by $t=0$
and then by $z=2$ on the residue curve.  Its final residue is $\Q$.
The iterated coefficientwise Puiseux construction with compatible
roots of $t$ and $z-2$ splits its decomposition sequence over
$G_\Q$, just as in B.2 at each step.  One may realize this
construction in iterated Laurent series, and restrict the action
to the separable closure of $\Q(t,z)$ therein.  Its composite with
$G_{\Q(t,z)}\twoheadrightarrow\pi_1(F)$ is a section, and
restriction to $G_{\Q(t)}$ is the coefficientwise cusp splitting
at $t=0$.  These are actual restriction and generic-point maps.

The marked curve over $\Q((t))$ is
$(\Pj^1;0,1,\infty,t)$.  Its ordered cross-ratio is $t$.
For a smooth four-marked model with disjoint sections, the
cross-ratio and its complement are units, since reduction is a
point of $M_{0,4}$ rather than its three boundary points.
Here $v(t)>0$, and that remains so after every finite extension
of the valued field.  Good marked reduction is consequently
impossible after every such extension.

For each one-variable intermediate field $L$, restrict the fixed
prolonged valuation by intersection with $L^{\mathrm{sep}}$ in
the chosen common closure.  Restriction preserves invariance under
the constructed section; intersections commute for nested fields.
The resulting valuation on $L$ has residue $\Q$, since its residue
embeds into the final residue $\Q$.  It is nontrivial, since a
trivial valuation would embed the transcendental field $L$ into
that residue.  On the smooth proper curve with field $L$ it is
therefore the valuation of a rational point.  This proves the
claimed compatible family for this example.  It supplies no
construction for an arbitrary section above a bad marked cusp.
\end{proof}

\begin{statement}{B.15}{Open}{The remaining birational and valuative assertions}
For an arbitrary continuous section $s$ of the tripod over $\Q$,
we have not split
\[
 1\longrightarrow\ker q_U\longrightarrow
 G_{\Q(t)}\times_{\pi_1(U),s}G_\Q
 \longrightarrow G_\Q\longrightarrow1.
\]
For an arbitrary section of $G_{\Q(t)}\to G_\Q$, we have not
proved that its image lies in a rational-place decomposition group
(B1, cote 149 scans 120--130).  B.1 gives local centers and density-one
integrality, not the missing global rational valuation.

For an arbitrary section of a hyperbolic curve fibration over a
hyperbolic curve with cuspidal base projection and genuinely bad
marked reduction, we have not constructed a birational lift, or
a single invariant valuation restricting compatibly to every
one-variable subfield (B2).  B.14 proves these properties only
for its displayed section.  B.10 refutes universal distinctness
of the inertia images in scan 128.  B.11 refutes the unrestricted
center converse in scan 120; the added elementary-hyperbolic
scope in its margin remains undecided here.  B.12 refutes an
unrestricted minimal-surface injectivity assertion in the stated
sense, while the existence assertion for appropriately minimal
strict hyperbolic compactifications remains undecided.
\end{statement}

\section{Local rings}
% Local-ring fragment. Insert in Part D after D.1--D.4.
\subsection{Faithfulness and the openness test for regular local rings}

Throughout this subsection $A,B$ are regular local $\Q$-algebras
essentially of finite type.  We use chosen geometric generic points.
Morphisms and ring homomorphisms are over $\Q$.  Write
$E_A=\pi_1(A)$, $E_B=\pi_1(B)$ and $p_A,p_B$ for their arithmetic
maps to $G_\Q$.  To retain the constants when these maps are not
surjective, a relative outer homomorphism means an equivalence
class of pairs
\[
 (\phi,\gamma),\qquad
 \phi:E_A\longrightarrow E_B\text{ continuous},\quad
 \gamma\in G_\Q,\quad
 p_B\phi=\operatorname{Int}(\gamma)p_A,
 \tag{D.5.0}
\]
where
\[
 (\phi,\gamma)\sim
       (\operatorname{Int}(h)\phi,p_B(h)\gamma),\qquad h\in E_B.
\]
This is the explicit base-point transporter in an outer map over
$\operatorname{Spec}\Q$: a path from the image of the source
geometric point to the target point induces $\phi$, and its
arithmetic projection is $\gamma$; changing the path gives exactly
the displayed equivalence.  We write
$\operatorname{Hom}^{\rm out}_{/G_\Q}(E_A,E_B)$ for these classes.
When $p_B$ is surjective, choose $h$ with $p_B(h)=\gamma^{-1}$.
This reduces the pair to one with $\gamma=1$, and the remaining
conjugations have $p_B(h)=1$.  Thus this definition recovers
strict maps over $G_\Q$ modulo target geometric conjugation in
that case.

The transporter also gives the correct map on algebraic constants
without a geometric connectedness assumption.  For number fields
$A=k$, $B=l$ embedded in $\bar\Q$, the condition on the pair is
$\gamma G_k\gamma^{-1}\subset G_l$, and $\phi$ is its forced
inclusion.  This condition is equivalent to $\gamma^{-1}l\subset k$.
The map from its class to $\gamma^{-1}|_l:l\to k$ is a bijection
onto $\operatorname{Hom}_\Q(l,k)$: its fibers are exactly the
left $G_l$-cosets used in the equivalence relation.  This checks
the constants and the direction of the transporter explicitly.
For a general ring its algebraic constant field gives the
corresponding open arithmetic image; every geometric kernel below
is the kernel of the indicated map to $G_\Q$.

\begin{statement}{D.5}{Proved}{Power-separatedness for finitely generated fields}
For every field $F$ finitely generated over $\Q$,
\[
                 \bigcap_{n\ge1}(F^\times)^n=\{1\}.          \tag{D.5.1}
\]
\end{statement}
\begin{proof}
Let $k$ be the algebraic closure of $\Q$ in $F$, a number field,
and choose a normal proper $k$-model of $F$.
Suppose $a=b_n^n$ for some $b_n\in F^\times$ for every $n$.
Every prime-divisor valuation of $a$ is divisible by every $n$
and thus zero; the same valuations of $b_n$ are zero.
Normality and properness imply $a,b_n\in k^\times$: a rational
function with no codimension-one poles is globally regular, and
global regular functions on this proper geometrically integral
model are constants.  The assertion is immediate at this step
as well when $F=k$.

At every finite prime of $k$, the equality $a=b_n^n$ again forces
the valuations of $a$ and all $b_n$ to vanish.  Thus these elements
are algebraic-integer units.  For every embedding $\sigma:k\to\C$,
\[
 |\sigma(b_n)|=|\sigma(a)|^{1/n}
       \le\max\{1,|\sigma(a)|\}.
\]
There are only finitely many algebraic integers of bounded degree
and bounded conjugates: the coefficients of their monic minimal
polynomials are bounded integers by the elementary symmetric
formulas.  Some $b_n=b_m$ therefore occurs for distinct $n,m$,
so this element, and consequently $a$, is a root of unity.
Let $w$ be the order of the finite group $\mu(k)$ and use
$a=b_w^w$.  The element $b_w$ is also a root of unity and
$b_w^w=1$.  Hence $a=1$.
\end{proof}

\begin{statement}{D.6}{Proved}{The arithmetic fundamental-group functor is faithful on these spectra}
Two morphisms $\operatorname{Spec}A\to\operatorname{Spec}B$
inducing the same relative outer homomorphism (D.5.0) are equal.
No dominance assumption is required for this assertion.
\end{statement}
\begin{proof}
Let the corresponding ring maps be $f,g:B\to A$.  For every
$n\ge1$, the coefficient module $M=\mu_n$ comes from $G_\Q$.
For a pair $(\phi,\gamma)$ the exact coefficient transport is
$M$ with action $p_B\phi\to M$ with action $p_A$,
$m\mapsto\gamma^{-1}m$.  In additive cocycle notation its
pullback is therefore
\[
                  c\longmapsto
       \bigl[a\longmapsto\gamma^{-1}c(\phi(a))\bigr].
 \tag{D.6.1}
\]
The equality $p_B\phi=\operatorname{Int}(\gamma)p_A$ verifies
that this is a cocycle with the required source action.  Replacing
the pair by $(\operatorname{Int}(h)\phi,p_B(h)\gamma)$ changes
(D.6.1) by the coboundary of
$-\gamma^{-1}p_B(h)^{-1}c(h)$, as follows directly from
$c(huh^{-1})=p_B(h)c(u)+(1-p_B(huh^{-1}))c(h)$.
Thus the induced cohomology pullback depends only on (D.5.0).
For an actual scheme morphism this is exactly its pullback of
the sheaf $\mu_n$: its base-point path acts on that sheaf by
the same $\gamma$.  Equality of the two relative outer maps
therefore gives equal pullbacks on $H^1(B,\mu_n)$.
The degree-one \'etale/fundamental-group
comparison and the Kummer sequence give, for every $u\in B^\times$,
\[
                  f(u)/g(u)\in(A^\times)^n
                  \qquad\text{for every }n.
\]
Apply D.5 to $\operatorname{Frac}(A)$ to obtain $f(u)=g(u)$.
For arbitrary $b\in B$, choose $c\in\Q$ such that $b-c$ is a
unit.  Such a $c$ exists because the residue of $b$ excludes at
most one rational number.  The maps agree on that unit and on
$c$, so they agree on $b$.  Thus $f=g$.
\end{proof}

\begin{statement}{D.7}{Proved}{An actual morphism has open image exactly when it is dominant}
For a morphism $f:\operatorname{Spec}A\to\operatorname{Spec}B$,
the image of $f_*:\pi_1(A)\to\pi_1(B)$ is open if and only if
the ring map $B\to A$ is injective.  Consequently there is an
injective map
\[
 \operatorname{Hom}^{\mathrm{dom}}_\Q
       (\operatorname{Spec}A,\operatorname{Spec}B)
 \lhook\joinrel\longrightarrow
 \operatorname{Hom}^{\mathrm{out,open}}_{/G_\Q}
       (\pi_1(A),\pi_1(B)).                       \tag{D.7.1}
\]
\end{statement}
\begin{proof}
Suppose first that $B\to A$ is injective, and set
$F=\operatorname{Frac}B\subset E=\operatorname{Frac}A$.
The algebraic closure $F'$ of $F$ in $E$ is finite over $F$.
With compatible separable closures, restriction $G_E\to G_F$
has image $G_{F'}$: it is closed by compactness and its fixed
field inside $F^{\mathrm{sep}}$ is $E\cap F^{\mathrm{sep}}=F'$.
It is therefore open.  Since $A,B$ are normal, the generic-point
maps $G_E\twoheadrightarrow\pi_1(A)$ and
$G_F\twoheadrightarrow\pi_1(B)$ are surjective.  The commutative
square of these four groups identifies the image of $f_*$ with
the image of $G_{F'}$ under the second surjection.  Surjective
homomorphisms of profinite groups are open, so $f_*$ has open image.

Conversely take $0\ne b\in\ker(B\to A)$.  Such an element lies
in the maximal ideal of $B$, since a unit cannot map to zero.
The element $u=1+b$ is a unit, and its image in $A$ is $1$.
It is nonconstant over $\Q$: a nonzero element algebraic over
$\Q$ has inverse given by its minimal polynomial and is a unit
in every $\Q$-algebra containing it; thus $b$ cannot be algebraic.

Let $k$ be the algebraic closure of $\Q$ in $F=\operatorname{Frac}B$.
Normality gives $k\subset B$, since elements of $k$ are integral
over $B$.  The chosen geometric component has field
$F\bar\Q$, regular over $\bar\Q$.  On a normal proper
$\bar\Q$-model, the nonconstant rational function $u$ has a
prime-divisor valuation $d\ne0$.  For every $n$, adjoining
$u^{1/n}$ has ramification index at least $n/\gcd(n,d)$ there.
But $u$ is a unit on the geometric local scheme, so this Kummer
cover is finite \'etale on that scheme.  Choosing compatible roots
gives a homomorphism
\[
 \theta:\Pi_B\longrightarrow\widehat\Z(1),\qquad
 \Pi_B=\ker(\pi_1(B)\to G_\Q),
\]
with infinite image: the orders of its finite Kummer images have
the displayed unbounded lower bounds.

Write $H=\operatorname{im}f_*$.  The pullback of $u$ is $1$, so
its Kummer homomorphism vanishes on $f_*(\Pi_A)$.
Moreover $H\cap\Pi_B=f_*(\Pi_A)$: the arithmetic compatibility
is $p_Bf_*=\operatorname{Int}(\gamma)p_A$, so any preimage of
an element with trivial arithmetic projection also has trivial
arithmetic projection.
Thus $H\cap\Pi_B\subset\ker\theta$.  If $H$ were open, this
would be an open subgroup of $\Pi_B$, forcing the image of
$\theta$ to be finite, a contradiction.  This proves the converse.
D.6 now gives injectivity of (D.7.1).
\end{proof}

\begin{statement}{D.8}{Open}{The exact fullness assertion for D5}
For these regular local rings, it remains to construct a dominant
morphism from an arbitrary relative outer homomorphism
$(\phi,\gamma)$ of (D.5.0) whose continuous homomorphism
$\phi:\pi_1(A)\to\pi_1(B)$ has open image.
D.5--D.7 prove faithfulness and the exact openness
test for actual morphisms; they do not prove surjectivity of
(D.7.1).
\end{statement}

This scope is the spectra of regular local rings essentially of
finite type specified in cote 149, scans 49--50.  For comparison,
a smooth ample hyperplane section of a proper hyperbolic polycurve
of dimension at least two induces a surjection on fundamental
groups by Lefschetz, while its inclusion is not dominant.  This
proper-variety example is given in
\cite[Example 3.13]{BIR_Carlson2026}.  It does not give a map
between the local spectra in (D.7.1).  The local-ring calculation
exhibits the extra finite \'etale covers explicitly: if a local
ring map kills $b\ne0$, all the covers $z^n=1+b$ become trivial
after pullback and their geometric images have unbounded order.

\section{Full class-two comparison}
% Insert as A.38 after the coefficient-by-coefficient cup theorem A.27.
\begin{statement}{A.38}{Proved}{Every ordinary local unit-equation pair lifts to the full profinite class-two quotient}
Let $\Pi=\pi_1(U_{\bar\Q})$ and let
$\Gamma_3\Pi=\overline{[\Pi,\overline{[\Pi,\Pi]}]}$ be its third
closed lower-central subgroup.  Put $E_2=E_U/\Gamma_3\Pi$ and
$E_{\rm ab}=E_U/\overline{[\Pi,\Pi]}$.
For every $(A,B)\in\mathcal G$, the continuous abelian section
of $E_{\rm ab}\to G_\Q$ given by its Kummer classes has a
continuous lift through
\[
                       E_2\twoheadrightarrow E_{\rm ab}.     \tag{A.38.1}
\]
This is the full profinite class-two quotient, with its actual
arithmetic action and all coefficient primes, including two.
\end{statement}
\begin{proof}
Choose the rational tangent $\overrightarrow{01}$ and the standard
geometric generators $x,y$ around $0,1$.  Let $R=\widehat\Z$ and
define the profinite group
\[
 H(R)=R^3,\qquad
 (a,b,c)(a',b',c')=(a+a',b+b',c+c'+ab').                         \tag{A.38.2}
\]
Set $X=(1,0,0)$, $Y=(0,1,0)$ and $Z=(0,0,1)$.
Then $[X,Y]=Z$ under the convention $[X,Y]=XYX^{-1}Y^{-1}$,
the center and commutator subgroup are $0\times0\times R$, and
\[
                     X^aY^bZ^d=(a,b,ab+d).                    \tag{A.38.3}
\]
We verify that this is the claimed geometric quotient, rather than
just an abstract group with the same graded modules.  The map
$x\mapsto X$, $y\mapsto Y$ from the free profinite group $\Pi$
is onto by (A.38.3), and kills $\Gamma_3\Pi$.
For any profinite class-two group generated topologically by $u,v$,
the formula
\[
 (a,b,c)\longmapsto u^a v^b[u,v]^{c-ab}                        \tag{A.38.4}
\]
is a continuous homomorphism.  Powers with exponents in
$\widehat\Z$ are defined on the closed procyclic subgroups, and
the collecting formula follows first for integer exponents and
then by continuity.  The commutator is central.  Formula (A.38.4)
applied to the images of $x,y$ in $\Pi/\Gamma_3\Pi$ is an inverse
to the preceding map.  Hence
\[
                         \Pi/\Gamma_3\Pi\simeq H(R).          \tag{A.38.5}
\]
This argument also identifies (A.38.5) with
$\prod_\ell H(\Z_\ell)$ by the coordinate projections of $R$.

The actual action induced by the chosen rational tangent is
\[
 g(x)=x^{\chi(g)},\qquad
 g(y)=f_g^{-1}y^{\chi(g)}f_g,
                       \quad f_g\in\overline{[\Pi,\Pi]}.
                                                                    \tag{A.38.6}
\]
The base point, formula and membership of $f_g$ in the closed
commutator subgroup are verified in
\cite[\S2, equation (2.1) and the following paragraph]{BIR_NakamuraSchneps2000}.
Modulo $\Gamma_3\Pi$, the image of $f_g$ is central.  Applying
(A.38.3) shows that (A.38.6) becomes exactly
\[
          g(a,b,c)=(\chi(g)a,\chi(g)b,\chi(g)^2c).             \tag{A.38.7}
\]
There is no division by two in (A.38.2)--(A.38.7).  Thus the
rational tangential splitting identifies the actual arithmetic
quotient with $E_2=H(R)\rtimes G_\Q$ with this action.
Its abelianization is $R(1)^2\rtimes G_\Q$.  The functions
$t,1-t$ identify the two geometric abelian coordinates with their
Kummer characters: their orders at the two corresponding cusps
are $(1,0)$ and $(0,1)$.  On the tangent
$\overrightarrow{01}$ both Kummer classes are zero, from their
leading coefficients $1$.  This fixes the arithmetic calibration
of the same two coordinates.

Choose continuous one-cocycles $a,b:G_\Q\to R(1)$ representing
$A,B$.  The associated actual abelian section is
$g\mapsto(a(g),b(g);g)$.  Its set-theoretic lift
$g\mapsto(a(g),b(g),0;g)$ to (A.38.2) has central defect
\[
                     c(g,h)=a(g)\chi(g)b(h),                  \tag{A.38.8}
\]
the usual continuous cup cocycle with values in $R(2)$.
The coefficient conclusions of A.27 yield zero for its image
in every $H^2_{\rm cts}(G_\Q,\Z_\ell(2))$.
To justify the full profinite inference explicitly, continuity
into a product is coordinatewise, so there is an identity of
cochain complexes
\[
 C^\bullet_{\rm cts}(G_\Q,R(2))
       =\prod_\ell C^\bullet_{\rm cts}(G_\Q,\Z_\ell(2)).      \tag{A.38.9}
\]
Its kernels and images are the respective products.  Equivalently,
choose a continuous primitive $h_\ell$ of $-c_\ell$ in every
factor.  Their product $h=(h_\ell)_\ell$ is continuous and
satisfies $\delta h=-c$.  Therefore
\[
                    g\longmapsto(a(g),b(g),h(g);g)            \tag{A.38.10}
\]
is a continuous section of $E_2\to G_\Q$ and projects to the
chosen abelian representative.  Changing the abelian cocycles
by a coboundary corresponds to conjugation by an element of
$H(R)$ lifting that geometric abelian element.  This proves
the assertion on the section classes with their specified
geometric conjugation as well.

The program \texttt{full\_cup\_class2\_certificate.py} checks
the collecting and action identities over several finite rings,
including powers of two, and the original-cup tame residue
formula used in the coefficient proof.  The existence of the
continuous primitive in (A.38.10) is furnished by A.27 and the
cochain identity (A.38.9), not by those finite checks.
\end{proof}

The remaining lifting map is the actual continuous surjection
$E_U\twoheadrightarrow E_2$, with kernel $\Gamma_3\Pi$.
A.38 constructs no section through that kernel.  In particular
the class-two obstruction in L11 vanishes on $\mathcal G$;
A.44 supplies nonrational elements of $\mathcal G$. Their possible
lifts through higher quotients remain undecided, with these precise
comparison maps and the class-three calculation in A.43.

\bibliographystyle{alpha}
\bibliography{../../notes/birational_references}
\end{document}
